% Ad M. Brutum 1.6 (ad-brutum-01-06)
% Source: https://raw.githubusercontent.com/PerseusDL/canonical-latinLit/master/data/phi0474/phi059/phi0474.phi059.perseus-lat1.xml
% Fetched by scripts/fetch_latin.py

% Dateline (Perseus): Scr. Romae xiv K. Iun. a. 711 (43).

\ciceroLetterOpener{Brutus Ciceroni salutem}

\ciceroSection{1} noli exspectare dum tibi gratias agam. iam pridem hoc ex nostra necessitudine quae ad summam benevolentiam pervenit sublatum esse debet. filius tuus a me abest; in Macedonia congrediemur. iussus est enim Ambracia ducere equites per Thessaliam. scripsi ad eum ut mihi Heracleam occurreret. Cum eum videro, quoniam nobis permittis, communiter constituemus de reditu eius ad petitionem aut commendationem honoris.

\ciceroSection{2} tibi Glycona, medicum Pansae, qui sororem Achilleos nostri in matrimonio habet, diligentissime commendo. audimus eum venisse in suspicionem Torquato de morte Pansae custodirique ut parricidam. nihil minus credendum est; quis enim maiorem calamitatem morte Pansae accepit? praeterea est modestus homo et frugi quem ne utilitas quidem videatur impulsura fuisse ad facinus. rogo te et quidem valde rogo (nam Achilleus noster non minus quam aequum est laborat) eripias eum ex custodia conservesque. hoc ego ad meum officium privatarum rerum aeque atque ullam aliam rem pertinere arbitror.

\ciceroSection{3} Cum has ad te scriberem litteras, a Satrio, legato C. Treboni, reddita est epistula mihi a Tillio et Deiotaro Dolabellam caesum fugatumque esse. Graecam epistulam tibi misi Cicereii cuiusdam ad Satrium missam.

\ciceroSection{4} Flavius noster de controversia quam habet cum Dyrrhachinis hereditariam sumpsit te iudicem. rogo te, Cicero, et Flavius rogat rem conficias. quin ei qui Flavium fecit heredem pecuniam debuerit civitas non est dubium; neque Dyrrhachini infitiantur sed sibi donatum aes alienum a Caesare dicunt. noli pati a necessariis tuis necessario meo iniuriam fieri. x iiii K. Iunias ex castris ad imam Candaviam.
